% femtolisp 笔记
% license Usr
% type Note

\subsection{构建}
\begin{itemize}
\item 如果编译了 Julia，那么 \verb`src/flisp/flisp` 就是 REPL。
\item julia 中， \verb`julia/src` 或 \verb`julia/src/flisp` 有许多 scm 文件和 lsp 文件（都在 \verb`flsp` 中）。
\item \verb`julia-parser.scm` 是语法解析，\verb`ast.scm` 是语法树工具
\item 
\end{itemize}


\subsection{类型}
\begin{itemize}
\item \href{https://code.google.com/archive/p/femtolisp/wikis/Manual.wiki}{语法手册（非官方）}
\item \verb`(typeof 1)` 是 \verb`fixnum`
\item \verb`(typeof 'abc)` 是 \verb`symbol`， 又如 \verb`'+`
\item \verb`'abc` 是 \verb`(quote abc)` 的语法糖
\item \verb`(typeof 1.23e4)` 是 \verb`double`
\item \verb`(typeof "abc")` 是 \verb`array byte`
\item \verb`(typeof #t)` 是 \verb`boolean`， 还有 \verb`#f`
\item \verb`(typeof '())` 是 \verb`null`
\item \verb`(typeof (cons 1 2))` 是 \verb`pair`
\item \verb`(car 变量)` 和 \verb`(cdr 变量)` 分别取第 1 和 2 个元素
\item \verb`(cadr 变量)` 等效于 \verb`(car (cdr 变量))`，以此类推。
\item \verb`(list 1 2 3)` 等效为 \verb`(cons 1 (cons 2 (cons 3 '())))`， 本质上是单链表。 类型仍然是 \verb`pair`。
\item \verb`(typeof (vector 1 2 3))` 是 \verb`vector`
\item \verb`(aref (vector 10 11 12) 0)` 获取第 0 个元素 \verb`10`
\end{itemize}

\subsection{计算器}
\begin{itemize}
\item \verb`(+ 1 2 3)` 加法， 另有 \verb`-, *, /`
\item \verb`(define a 1)` 定义变量
\item \verb`(print 变量)` 打印 \verb`princ` 也可以。
\item \verb`(define a 123)` 然后 \verb|`(1 ,a)| 会得到 \verb`(1 123)`。
\item \verb`(define a '(1 2 3))` 然后 \verb|`(0 ,@a 4)| 得到 \verb`(0 1 2 3 4)`。
\item \verb`(cons 1 2)` 显示为 \verb`(1 . 2)`， \verb`(cons 1 (cons 2 (cons 3 4)))` 显示为 \verb`(1 2 3 . 4)`
\item \verb`(list 1 2 3 4)` 和 \verb`'(1 2 3 4)` 等效，显示为 \verb`(1 2 3 4)`
\item \verb`(eq 1 1)` 检查是否相等， 返回 \verb`#t`
\item \verb`(= 3 3)` 等效吗？
\item \verb`(length 列表)` 返回长度， \verb`(length= 列表 长度)` 检查长度是否符合预期。
\end{itemize}

\subsection{控制流}
\begin{itemize}
\item \verb`(if (> 2 1) 'big 'small)` 返回 \verb`'big`
\item \verb`(cond  ((= x 1) 'one)  ((= x 2) 'two) ((= x 3) 'three)  (else 'big))`
\end{itemize}

\subsection{字符串}
\begin{itemize}
\item \verb`(string "abc" "def" "ghi")` 拼接字符串
\item \verb`(string.join '("abc" "def") ", ")` 得到 \verb`"abc, def"`
\item \verb`(symbol "abc")` 把字符串转为符号。
\end{itemize}


\subsection{函数}
\begin{itemize}
\item 多次定义同名函数会把之前的定义覆盖。
\item \verb`(define (square n) (* n n))` 定义函数 \verb`(square 3)` 调用
\item \verb`(null? '())` 返回 \verb`#t`
\item \verb`(define (函数名 形参1 (形参2 默认值) ...) 函数体)`。 如 \verb`(define (myfun1 x (y 3.14)) (+ x y))`
\item 返回 \verb`boolean` 的函数称为 predicate，习惯上以 \verb`?` 结尾
\item 内建 predicate（以下全返回 \verb`#t`）：\verb`(number? 42)`，\verb`(odd? 3)`， \verb`(string? "abc")`， \verb`(null? '())`，\verb`(pair? '(1 2))`，\verb`(even? 6)`，\verb`(eq? 1 1)`
\item \verb`(map 函数 列表)` 会把函数依次作用在列表的每个元素并返回一个新的列表。 例如 \verb`(map abs '(-1 -2 3 -4))` 得到 \verb`(1 2 3 4)`。
\item \verb`*名字*` 只是个普通名字，习惯上用这种格式表示它是一个全局变量（声明在任何函数体之外）。
\end{itemize}

\subsection{循环}
C++ 中的
\begin{lstlisting}[language=cpp]
int f() {
    int a = 1, b = 2, c = 3;
    for (int i = 0; i < 100; i++) {
        a = a + b;
        b = b + c;
        c = c + a;
    }
    return a + b + c;
}
\end{lstlisting}
翻译为
\begin{lstlisting}[language=none]
(define (f i a b c)
  (if (>= i 100)
      (+ a b c)
      (f (+ i 1)
         (+ a b)
         (+ b c)
         (+ c a))))
\end{lstlisting}
